\subsection{Übersicht}

Mit der Beschreibung in dem vorherigen Abschnitt haben wir einige relevante Punkte der Entwicklung bereits angeschnitten.

Im Folgenden geben wir einen Überblick über die Architektur sowie der Kernkonzepte, die in helios zur Konfiguration, Kommunikation und Ressourcenverwaltung der in der GameWorld aktiven GameObjects und Systeme implementiert wurden.

\subection{Strukturierung der engine-Schicht}

Die Strukturierung der engine Schicht folgt der im ersten Dokument beschriebenen Unterteilung in Module und folgt wie in dem dem Dokument weitestgehend einer ``by-feature, by-layer``-Strukturierung und ist in Abbildung~\ref{fig:engine_struktur} dargestellt.

\begin{figure}[htbp]
    \setlength{\DTbaselineskip}{18pt}
    \dirtree{%
        .1 ./engine
        .2 builder/.
        .2 common/.
        .2 ecs/.
        .2 mechanics/.
        .2 modules/.
        .2 runtime/.
        .2 state/.
        .2 tooling/.
    }
    \caption{Struktur des engine-Moduls von helios.}
    \label{fig:engine_struktur}
\end{figure}

Hierbei gestaltet sich die Unterteilung in Module wie folgt:

\begin{itemize}
    \itemsep0.5em
    \item \texttt{builder/}: Stellt Builder-Klassen für Spawn-Systeme und GameObjects-bereit, auf Basis von Fluent Interface sowie einer Domain-Specific Language (DSL), um die Erstellung von Fachobjekten von einem hohen Abstraktionsgrad in einen zugänglichen semantischen Kontext zu überführen.
    \item \texttt{common/}: Enthält allgemeine Hilfsklassen, die vor allem von \texttt{mechanics} und \texttt{modules} genutzt werden, darunter spielspezifische Kontextklassen, die für verschiedene Systeme und Komponenten relevante Informationen über die aktuelle Spielsituation bereitstellen.
    \item \texttt{core/}: Stellt fundamentale Datentypen und -Objekte zur Verfügung, die für den Betrieb der engine benötigt werden, darunter bspw. \texttt{Units}, die das metrische System für die helios bereitstellt.
    \item \texttt{ecs/}: Implementiert die Kernkomponenten der ECS-Architektur, darunter den \texttt{EntityManager}, die \texttt{GameObject}-Fassade und die \texttt{SparseSet}-Datenstruktur.
    \item \texttt{mechanics/}: Implementiert Komponenten, Systeme, Datentypen und Nachrichten-Objekte für verschiedene Spielmechaniken, beispielsweise für Eingabeverarbeitung, Gesundheits- und Schadenszustände sowie Punktesysteme.
    \item
\end{itemize}